Mountain ridges and river valleys are among the oldest natural boundaries in human settlement — communities on opposite sides of a ridge typically have different economies, cultures, and administrative histories. This paper operationalises topographic separation as an edge weight modifier for the \textsc{bisect} redistricting pipeline. Using USGS 3DEP 1/3 arc-second digital elevation model (DEM) data at 10m resolution, we derive three topographic signals for each census-tract adjacency edge: ridge-crossing penalty (reducing weight to 0.1$\times$ when an edge crosses a watershed divide), elevation-similarity bonus (exponential decay over 300m elevation difference), and valley-clustering bonus (1.5$\times$ weight for tracts sharing a watershed below the ridge). Applied to North Carolina's five mountain congressional districts ($k=5$), topographic weights produce zero cross-ridge districts — compared to 2--3 cross-ridge districts under standard geographic weights — while maintaining Polsby-Popper compactness within 4\% of the unweighted baseline. For flat states (Wisconsin plains, Texas coastal plain) where elevation variation is $<50$m across all tracts, topographic weights degrade gracefully to near-baseline behavior: the similarity metric is uniformly high and the edge weights are unmodified. We provide the topographic weight pipeline as the \texttt{--weights-override topographic} flag in the \textsc{bisect} CLI, using only public-domain USGS data and no proprietary datasets. Legal grounding is the communities-of-interest doctrine, which courts have interpreted to include shared geography, shared economic activity, and shared administrative membership — all of which are captured by watershed-level topographic clustering.
